% !TEX root = ../LAFROX-Formulario-T-13.tex
\chapter*{Formulario T-13: Plan de pruebas y validación}
\addcontentsline{toc}{chapter}{Formulario T-13: Plan de pruebas y validación}
\markboth{Formulario T-13: Plan de pruebas y validación}{}
\label{cap:t13}

Este formulario acompaña al Subdocumento 9, Plan de calidad (en adelante, SD9). Contiene el plan de pruebas y validación que exige el Formulario T-13 de las Bases Administrativas: niveles, tipos, ambientes, datos de prueba, criterios de entrada y salida, automatización, y calendario de las pruebas de carga, resiliencia, recuperación ante desastres y seguridad ofensiva (Distribuidora Puelche S.A., 2026a, Formulario T-13, p. 62). El SD9 resume y analiza este plan; el protocolo con el que el CLIENTE acepta cada hito está en el Formulario T-17.

En este formulario, BA, BTT y Caso designan las Bases Administrativas (Distribuidora Puelche S.A., 2026a), las Bases Técnicas Transversales (2026b) y las Bases Técnicas del Caso Logística (2026c). Mes 1 del contrato = febrero de 2027. Los paquetes de trabajo son los del Formulario T-14 y las fechas y horas, las del Formulario T-15. CAL, SEG, SRE, DAT e IMP designan los roles de calidad y pruebas, seguridad, operación y confiabilidad, datos e implantación del Formulario T-15; la sección 11 indica quién los ejerce.

\section*{1. Objeto, alcance y fuentes}
\addcontentsline{toc}{section}{1. Objeto, alcance y fuentes}

El plan cubre las pruebas de la implementación (meses 1 a 21), de las dos marchas blancas (meses 13 a 15 y 19 a 20) y de la Operación (meses 21 a 56). Se aplica al monolito modular de 12 módulos (M1 a M12), a las 15 integraciones (INT-01 a INT-15), a la aplicación móvil Android en sus cuatro perfiles, a las consolas y portales web, a la infraestructura híbrida y a la migración de datos, tal como los define el Subdocumento 4 (SD4).

Las exigencias que el plan satisface son las siguientes:

\begin{tablasimple}{>{\hsize=0.559\hsize\linewidth=\hsize}Y >{\hsize=1.441\hsize\linewidth=\hsize}Y}
  {\cab{Fuente} & \cab{Exigencia que rige este plan}}
Bases Administrativas, Art. 4.3 (BA, pp. 5--6) & ISO/IEC/IEEE 29119 para pruebas; ISO/IEC 25010 y 25012 para calidad del producto y de los datos \\
Bases Administrativas, Art. 24 (BA, p. 17) & Cuatro ambientes aislados; regresión automatizada; cobertura ≥ 70 \% de la lógica de negocio con umbral bloqueante; carga y estrés a 1,5 × el peak; resiliencia antes de cada paso a producción y semestral \\
Bases Técnicas Transversales (BTT), sección 20.1 (BTT, pp. 34--35) & Diez tipos de prueba, con su momento y su criterio de salida \\
BTT, cap. 4 (BTT, pp. 10--11) & Cinco ambientes y sus exigencias de datos (4.1); trazabilidad requisito–incidencia–código–prueba–despliegue (RT-04.04); flujo de integración con bloqueo automático (RT-04.05) \\
BTT, cap. 9 (BTT, pp. 21--22) & Pruebas de carga y de estrés hasta el punto de quiebre (RT-09.06), con el contenido mínimo de su informe (RT-09.07), y degradación controlada (RT-09.08) \\
BTT, RT-07.07, RT-07.12, RT-10.07, RT-11.20 (BTT, pp. 17--18, 22, 24) & Pruebas periódicas de la Operación \\
Caso Logística, cap. 15 & Umbrales de respuesta (RT-09.01), operación desconectada (RT-03.10, RT-03.13), ventana protegida y congelamientos (RT-10.05), interfaces de terreno (RT-13.08) \\
Caso Logística, cap. 18 & Dieciséis resultados de aceptación del caso (R18-01 a R18-16) \\
\end{tablasimple}
\fuentesd{Fuente: elaboración a partir de Distribuidora Puelche S.A. (2026a, arts. 4.3 y 24, pp. 5--6, 17); Distribuidora Puelche S.A. (2026b, caps. 4, 7, 9, 10, 11 y 20, pp. 10--11, 17--18, 21--24, 34--35); Distribuidora Puelche S.A. (2026c, caps. 15 y 18).}

\section*{2. Proceso de pruebas según ISO/IEC/IEEE 29119}
\addcontentsline{toc}{section}{2. Proceso de pruebas según ISO/IEC/IEEE 29119}

El plan sigue los tres niveles de proceso de ISO/IEC/IEEE 29119-2 (ISO, IEC \& IEEE, 2021a):

\begin{tablasimple}{>{\hsize=0.478\hsize\linewidth=\hsize}Y >{\hsize=1.643\hsize\linewidth=\hsize}Y >{\hsize=1.245\hsize\linewidth=\hsize}Y >{\hsize=0.634\hsize\linewidth=\hsize}Y}
  {\cab{Nivel de proceso} & \cab{Artefacto} & \cab{Responsable} & \cab{Revisión}}
Organizacional & Política de pruebas de LafroX, dentro del sistema de gestión de calidad ISO 9001:2015 & Comité de Calidad y Procesos de LafroX & Anual \\
Gestión de pruebas & Este plan (T-13) y su seguimiento por campaña & Líder de Calidad y Pruebas, Maximiliano Miño & En cada hito y ante cada cambio aprobado de alcance \\
Pruebas dinámicas & Diseño, preparación del ambiente, ejecución e informe de cada caso & Equipo de pruebas (rol CAL) y especialistas (SEG, SRE, DAT, IMP) & En cada ejecución \\
\end{tablasimple}
\fuentesd{Fuente: elaboración a partir de ISO, IEC \& IEEE (2021a, procesos de prueba); LafroX (2026i, SD1, secciones 1.3.1 y 1.5); LafroX (2026o, SD6, organización del desarrollo).}

La intensidad de la prueba sigue el riesgo del producto, como dispone ISO/IEC/IEEE 29119-1 (ISO, IEC \& IEEE, 2022). La prioridad de cada requisito en el catálogo del SD3 (crítica, alta o media) fija cuántos casos se diseñan para él (sección 8). La documentación de pruebas sigue la estructura de ISO/IEC/IEEE 29119-3 (ISO, IEC \& IEEE, 2021c): especificación de casos, procedimientos, registro de ejecución, registro de incidencias e informe de finalización.

\subsection*{2.1 Técnicas de diseño (ISO/IEC/IEEE 29119-4)}
\addcontentsline{toc}{subsection}{2.1 Técnicas de diseño (ISO/IEC/IEEE 29119-4)}

El equipo elige la técnica según la regla que debe comprobar. La tabla muestra su aplicación a los procesos de Puelche.

\begin{tablasimple}{>{\hsize=0.533\hsize\linewidth=\hsize}Y >{\hsize=0.659\hsize\linewidth=\hsize}Y >{\hsize=1.808\hsize\linewidth=\hsize}Y}
  {\cab{Tipo de regla} & \cab{Técnica} & \cab{Ejemplo en Puelche}}
Umbrales numéricos & Partición de equivalencia y valores límite & Rango térmico de un producto refrigerado; límite de crédito del cliente en la preventa \\
Reglas con varias condiciones & Tablas de decisión & Asignación de stock por vencimiento y lote; entrega con local cerrado \\
Ciclos de vida & Transición de estados & Estados del pedido, de la guía de despacho y de la misión de preparación \\
Flujos de usuario & Casos de uso y escenarios & Flujos del SD3 por perfil: preventista, preparador, conductor, bodega \\
Integraciones & Pruebas de contrato & Siete casos por integración: éxito, sin respuesta, con error, con lentitud, duplicado, orden y versión (RT-10.08) (BTT, p. 22) \\
Concurrencia y reconexión & Escenarios con conflicto & Dos dispositivos que reservan el mismo lote sin señal y sincronizan después \\
\end{tablasimple}
\fuentesd{Fuente: elaboración a partir de ISO, IEC \& IEEE (2021b, técnicas de diseño); Distribuidora Puelche S.A. (2026c, caps. 13 y 15); Distribuidora Puelche S.A. (2026b, RT-10.08, p. 22).}

\section*{3. Niveles y tipos de prueba}
\addcontentsline{toc}{section}{3. Niveles y tipos de prueba}

Cada nivel y tipo se vincula con su ambiente, herramientas, grado de automatización y paquete de trabajo; esa correspondencia determina quién prepara y ejecuta la campaña.

\begin{tablasimple}{>{\hsize=1.084\hsize\linewidth=\hsize}Y >{\hsize=0.891\hsize\linewidth=\hsize}Y >{\hsize=2.231\hsize\linewidth=\hsize}Y >{\hsize=0.59\hsize\linewidth=\hsize}Y >{\hsize=0.558\hsize\linewidth=\hsize}Y >{\hsize=0.646\hsize\linewidth=\hsize}Y}
  {\cab{Nivel o tipo} & \cab{Ambiente} & \cab{Herramienta} & \cab{Automatización} & \cab{Paquete T-14} & \cab{Exigencia}}
Unitaria y de componente & Desarrollo y cadena de integración & PHPUnit con PCOV (PHP); JUnit 5 con Kover (Kotlin); Jest (TypeScript) & Total & 1.5.2 & Art. 24; RT-04.11 (BA, p. 17; BTT, p. 11) \\
Contrato de API y eventos & Cadena de integración & Pact; Spectral contra OpenAPI 3.1 y AsyncAPI 2.6 & Total & 1.5.2 & Art. 4.3; RT-10.08 (BA, pp. 5--6; BTT, p. 22) \\
Integración de los 15 flujos & QA & Pruebas automatizadas con simuladores del ERP y de terceros & Total & 3.8.1; 3.9.1 & BTT 20.1 (BTT, pp. 34--35) \\
Sistema y regresión & QA & Playwright (consolas y portales); Espresso (Android) & ≥ 80 \% de los casos & 3.8.1; 3.9.1 & BTT 20.1; Art. 24 (BA, p. 17; BTT, pp. 34--35) \\
Análisis estático y de seguridad & Cadena de integración & PHPStan con Larastan; PHPMD, PHPCPD, PhpMetrics, deptrac; detekt; ESLint; SAST, SCA, secretos e imágenes & Total & 1.5.2 & RT-04.05 (BTT, p. 10) \\
Análisis dinámico de seguridad & QA & OWASP ZAP (web y API); MobSF (Android) & Total & 3.8.1; 3.9.1 & Art. 4.3 (OWASP) (BA, pp. 5--6) \\
Accesibilidad & QA y Preproducción & axe-core y revisión manual con teclado y lector de pantalla & Parcial & 3.8.2; 3.9.2 & RT-13.01 (BTT, p. 26) \\
Usabilidad con personas usuarias & Prototipos y Preproducción & Pruebas moderadas por perfil & Manual & 2.6.2; 2.6.3; 3.8.2; 3.9.2 & RT-13.03; RT-13.04 (BTT, p. 26) \\
Aceptación de usuario & Preproducción & Casos del SD3 ejecutados por personas del CLIENTE & Manual & 3.8.2; 3.9.2 & BTT 20.1 (BTT, pp. 34--35) \\
Perfil operacional de terreno & Terreno, cámara de congelado y Preproducción & Dispositivos reales, Espresso para los flujos repetibles & Parcial & 3.8.3; 3.9.1; 3.9.2 & RT-03.10; RT-03.13; RT-13.08 \\
Carga y estrés & Preproducción & k6 & Total & 3.8.4; 3.9.3 & Art. 24; RT-09.06 a RT-09.08 (BA, p. 17; BTT, pp. 21--22) \\
Resiliencia & Preproducción & AWS Fault Injection Service; desconexiones controladas del sitio emulado; corte de gateways de frío y de enlaces & Total en la inyección; manual en la observación & 3.8.4; 3.9.3; 8.2.2 & RT-10.07 (BTT, p. 22) \\
Continuidad de la arquitectura híbrida & Preproducción y sitio emulado & Nueve protocolos del SD4, Anexo 4-V (sección 7.5) & Parcial & 3.8.1; 3.8.3 a 3.8.5; 3.9.1; 3.9.3; 3.9.4 & RT-03.10; RT-07.04 (BTT, p. 17) \\
Recuperación ante desastres & Recuperación ante Desastres (us-east-1) & Conmutación regional real y recuperación del WMS de Talca en la nube (AL-DR-01) & Procedimiento automatizado con verificación manual & 3.8.5; 3.9.4; 8.1.3 & RT-07.07 (BTT, p. 17) \\
Restauración de respaldos & Producción y Recuperación & Restauración de una muestra representativa & Procedimiento automatizado & 4.2.1; 4.2.2; 4.3.1; 8.1.6 & RT-07.12 (BTT, p. 18) \\
Seguridad ofensiva & Preproducción con versión congelada & Tercero independiente & Manual & 3.8.6; 3.9.5; 8.2.2 & RT-11.20 (BTT, p. 24) \\
Despliegue sin interrupción y reversión & Preproducción & Cadena de entrega del SD4 (canario con reversión automática) & Total & 3.8.8; 3.9.7 & Art. 24; RT-04.07 (BA, p. 17; BTT, p. 11) \\
Migración de datos & Preproducción & Procedimiento de extracción, carga y conciliación del SD5 & Parcial & 3.7.1 a 3.7.5 & BTT 20.1; RT-05.11 (BTT, pp. 12, 34--35) \\
\end{tablasimple}
\fuentesd{Fuente: elaboración a partir de Distribuidora Puelche S.A. (2026a, arts. 4.3 y 24, pp. 5--6, 17); Distribuidora Puelche S.A. (2026b, secciones 4.1 y 20.1, pp. 10--11, 34--35); Distribuidora Puelche S.A. (2026c, cap. 15); LafroX (2026m, SD4, sección 4.2.4.1); LafroX (2026a, T-10, apartado DevSecOps); LafroX (2026d, T-14, paquetes indicados en cada fila).}

Las herramientas son las declaradas en el Formulario T-10 (SD6) y en la cadena de entrega del SD4. Este plan no agrega herramientas nuevas.

\section*{4. Ambientes de prueba}
\addcontentsline{toc}{section}{4. Ambientes de prueba}

Los cinco ambientes se diferencian por sus datos y por el control que permiten. Antes de cada ciclo se verifica el estado indicado en la tabla.

\begin{tablasimple}{>{\hsize=0.33\hsize\linewidth=\hsize}Y >{\hsize=1.777\hsize\linewidth=\hsize}Y >{\hsize=0.745\hsize\linewidth=\hsize}Y >{\hsize=1.148\hsize\linewidth=\hsize}Y}
  {\cab{Ambiente} & \cab{Pruebas que aloja} & \cab{Datos} & \cab{Estado antes de cada ciclo}}
Desarrollo & Unitarias, de componente y de contrato & Sintéticos & Reconstruible desde código \\
QA & Integración, sistema, regresión, análisis dinámico, accesibilidad automatizada, carga de regresión semanal & Juegos sintéticos versionados & Restituido a un estado conocido (BTT 4.1, p. 10) \\
Preproducción & Aceptación, carga y estrés, resiliencia, ensayo del paso a producción, migración, seguridad ofensiva & Sintéticos con volumen de peak; históricos anonimizados solo para migración & Equivalente a Producción en topología, configuración y versiones (RT-04.02) (BTT, p. 10) \\
Producción & Verificación posterior al despliegue, indicadores diarios de la marcha blanca, restauración de respaldos & Reales & No se prueba con datos ficticios \\
Recuperación ante Desastres & Conmutación real & Réplica productiva protegida & No se usa como ambiente funcional de prueba \\
\end{tablasimple}
\fuentesd{Fuente: elaboración a partir de Distribuidora Puelche S.A. (2026b, sección 4.1, p. 10); LafroX (2026m, SD4, sección 4.2.4.1); LafroX (2026n, SD5, sección 5.3).}

\subsection*{4.1 Instancias de Preproducción por campaña}
\addcontentsline{toc}{subsection}{4.1 Instancias de Preproducción por campaña}

Varias campañas coinciden en el calendario, pero las que requieren estados incompatibles no comparten la misma instancia al mismo tiempo. La resiliencia necesita carga sostenida; la prueba de intrusión exige una versión congelada; la migración carga datos históricos; la aceptación necesita datos estables. Por eso migración, intrusión y aceptación usan instancias aisladas; carga, resiliencia y despliegue usan sesiones sucesivas en PRE-C. Cada instancia de Preproducción se reconstruye desde el mismo código y con la misma topología y configuración. Siguen siendo Preproducción y no constituyen un sexto ambiente; se destruyen al cerrar la campaña (RT-04.02; RT-04.14) (BTT, pp. 10--11).

\begin{tablasimple}{>{\hsize=0.338\hsize\linewidth=\hsize}Y >{\hsize=0.973\hsize\linewidth=\hsize}Y >{\hsize=1.689\hsize\linewidth=\hsize}Y}
  {\cab{Instancia} & \cab{Campaña} & \cab{Condición que exige separarla}}
PRE-M & Ensayos de migración & Carga datos históricos anonimizados, que se eliminan al cerrar cada ensayo \\
PRE-C & Carga y estrés, luego resiliencia bajo carga, luego despliegue sin interrupción & La resiliencia y el despliegue se miden bajo la carga sostenida, en secuencia \\
PRE-S & Seguridad ofensiva & Versión congelada; no admite inyección de fallas en paralelo \\
PRE-A & Aceptación de usuario y accesibilidad & Las personas usuarias necesitan datos estables \\
\end{tablasimple}
\fuentesd{Fuente: elaboración a partir de Distribuidora Puelche S.A. (2026b, RT-04.02, p. 10); LafroX (2026m, SD4, sección 4.2.4.1); LafroX (2026n, SD5, sección 5.3); las condiciones de aislamiento de campañas descritas en esta sección 4.1.}

\textbf{Secuencia de sesiones en PRE-C.} Carga y resiliencia forman una campaña compatible; despliegue y reversión usan después el mismo entorno, no se ejecutan sobre él mientras se inyectan fallas. La preparación de scripts, revisión y elaboración de informes pueden correr en paralelo fuera de PRE-C. Las subventanas siguientes se incluyen en las fechas y HH de los paquetes existentes y no añaden esfuerzo a la línea base.

\begin{tablasimple}{Y Y Y Y}
  {\cab{Etapa y paquetes} & \cab{Sesiones de carga y resiliencia} & \cab{Restitución y comprobación} & \cab{Sesiones de despliegue y reversión}}
E1: 3.8.4 y 3.8.8 & 21–28 octubre 2027 & 29 octubre; snapshot de datos y versión, reinicio de colas y comprobación de observabilidad & 2–5 noviembre; análisis e informes de ambos paquetes hasta 9 noviembre \\
E2: 3.9.3 y 3.9.7 & 19–22 mayo 2028; el 21 de mayo no se programa trabajo humano ordinario & 23 mayo; mismo procedimiento de restitución & 24–25 mayo; informes hasta 26 mayo \\
\end{tablasimple}
\fuentesd{Fuente: elaboración a partir de LafroX (2026e, T-15, sección 6.1, paquetes 3.8.4, 3.8.8, 3.9.3 y 3.9.7); la separación de sesiones y restitución definida en esta sección 4.1.}

La prueba continua de aislamiento de 24 horas requiere una ventana y un turno explícitos; si cruza un día no laborable, se autoriza el turno de observación antes de empezar. La fecha de sesión no implica dotación disponible. CAL y SRE confirman personas nominales y capacidad neta, descontando recuperación, soporte y otros frentes, antes de aprobar la campaña. Las 320/320 HH de E1 y 320/160 HH de E2 son esfuerzo total por paquete, no duración de las sesiones invasivas ni cuatro campañas que puedan ocupar simultáneamente PRE-C. Si la dotación vigente no alcanza, el Jefe de Proyecto debe asignar recursos y aprobar la nivelación antes de G5; no se presume suficiencia del escenario histórico del T-15.

Sólo una sesión puede reservar PRE-C a la vez. Su salida deja un informe con el digest, estado de datos, colas y controles; la siguiente sesión no inicia sin restitución aprobada y carga basal dentro de umbral. PRE-M, PRE-S y PRE-A permanecen aisladas. El Arquitecto confirma capacidad, cuotas y equivalencia del conjunto de instancias antes de la primera campaña; los resultados se registran en el plan del H1 y se vuelven a comprobar al preparar H5/H10.

\section*{5. Datos de prueba}
\addcontentsline{toc}{section}{5. Datos de prueba}

\textbf{Regla.} Ningún ambiente no productivo contiene datos productivos reales sin anonimización verificada. La regla protege los datos personales de clientes, conductores de terceros y la geolocalización de trabajadores (Ley N.º 21.719, 2024) y los campos que el Caso, RT-11.10, exige cifrar: comportamiento de pago, antecedentes comerciales y geolocalización.

\begin{tablasimple}{>{\hsize=0.615\hsize\linewidth=\hsize}Y >{\hsize=1.022\hsize\linewidth=\hsize}Y >{\hsize=1.581\hsize\linewidth=\hsize}Y >{\hsize=0.422\hsize\linewidth=\hsize}Y >{\hsize=1.36\hsize\linewidth=\hsize}Y}
  {\cab{Juego} & \cab{Origen} & \cab{Volumen} & \cab{Ambiente} & \cab{Control}}
Maestros sintéticos & Generador con las reglas del SD3 y del SD5 & 14.200 clientes, 8.400 SKU (1.100 refrigerados o congelados), 180 proveedores (Caso, 14.1) & QA y Preproducción & Versionado; restituible \\
Transacciones de régimen & Generador sintético & 31.000 pedidos y 260.000 líneas de preparación al mes (Caso, 14.1) & QA y Preproducción & Versionado \\
Transacciones de peak & Generador sintético & Perfil de septiembre: \ensuremath{\approx} 2.600 entregas por día (Caso, B.2) & Preproducción (PRE-C) & Versionado \\
Lotes y temperatura & Generador sintético & Lotes con vencimiento; series de temperatura con excursiones controladas & QA y Preproducción & Incluye los casos límite del rango térmico \\
Históricos para migración & Datos del ERP, del WMS de 2013 y de planillas del CLIENTE & 32,11 GB (SD4, sección 4.2.6, dimensión 10) & Solo PRE-M & Anonimización verificada con Amazon Macie (0 hallazgos) más una prueba de reidentificación sobre una muestra; se eliminan al cerrar el ensayo \\
Incidentes reproducidos & Evidencia de terreno capturada por INN-02, sin datos personales & Según incidentes & Ambientes efímeros por incidencia en QA o Preproducción & Convertidos en casos de regresión permanentes (SD13, sección 13.2) \\
\end{tablasimple}
\fuentesd{Fuente: elaboración a partir de Distribuidora Puelche S.A. (2026c, sección 14.1 y cap. B); LafroX (2026n, SD5, sección 5.3); LafroX (2026m, SD4, secciones 4.2.6.3 a 4.2.6.5).}

\section*{6. Criterios de entrada y de salida}
\addcontentsline{toc}{section}{6. Criterios de entrada y de salida}

La campaña comienza sólo cuando cumple sus condiciones de entrada y termina con evidencia de los criterios de salida; la suspensión no equivale a aprobación.

\subsection*{6.1 Por nivel de prueba}
\addcontentsline{toc}{subsection}{6.1 Por nivel de prueba}

Las condiciones siguientes se verifican antes y después de cada ejecución. Las fallas bloqueantes se resuelven antes de promover el artefacto.

\begin{tablasimple}{>{\hsize=0.332\hsize\linewidth=\hsize}Y >{\hsize=1.308\hsize\linewidth=\hsize}Y >{\hsize=1.36\hsize\linewidth=\hsize}Y}
  {\cab{Nivel} & \cab{Criterio de entrada} & \cab{Criterio de salida}}
Unitaria y componente & Cambio con solicitud de fusión que cita su requisito o incidencia & Cobertura de lógica de negocio ≥ 70 \% y global ≥ 80 \%; 0 pruebas en falla; complejidad ciclomática ≤ 10 por método; 0 dependencias prohibidas entre módulos; acoplamiento eferente ≤ 20 clases por clase; duplicación ≤ 3 \% (T-10) \\
Contrato & Contratos OpenAPI 3.1 y AsyncAPI 2.6 versionados & 0 contratos rotos \\
Integración & Imagen promovida a QA sin vulnerabilidades críticas ni altas; simuladores de los 15 flujos disponibles & 15 de 15 flujos sin error (BTT 20.1, pp. 34-35) \\
Sistema y regresión & Integración aprobada; juego de datos restituido & 100 \% de la batería aprobada; 0 regresiones; p95, errores y colas de la carga de regresión no peores que en la última versión aceptada (SD4, sección 4.2.4.1.4) \\
Análisis dinámico & Versión desplegada en QA & 0 alertas altas de DAST; 0 hallazgos críticos o altos de MobSF \\
Aceptación de usuario & Regresión aprobada; 0 defectos críticos o altos abiertos; personas del CLIENTE convocadas y con inducción & Casos de aceptación firmados por la Contraparte Técnica (BTT 20.1, pp. 34-35) \\
Accesibilidad & Vistas completas en QA & 0 incumplimientos WCAG 2.2 A y AA; informe de conformidad entregado (RT-13.01) (BTT, p. 26) \\
Carga & Versión candidata en PRE-C; juego de peak cargado; monitoreo activo & Todos los p95 de la sección 7.1 a 21,99 TPS; 0 transacciones perdidas \\
Estrés & Carga aprobada & Punto de quiebre identificado; degradación con encolamiento y mensaje explícito, sin pérdida (RT-09.08) (BTT, p. 22) \\
Resiliencia & Carga sostenida en curso; escenarios de falla aprobados por el Arquitecto de Solución & Recuperación sin intervención en cada escenario; 0 pérdidas (RT-10.07) (BTT, p. 22) \\
Recuperación ante desastres & Procedimiento de conmutación aprobado; conmutación fuera de congelamiento, de la ventana de 05:30 a 07:00 y de los tres primeros días hábiles del mes & RTO ≤ 4 h y RPO ≤ 15 min medidos en conmutación real, regional y del WMS de Talca (RT-07.04, RT-07.07; AL-DR-01) (BTT, p. 17) \\
Seguridad ofensiva & Versión congelada en PRE-S; alcance acordado con el tercero & 0 hallazgos críticos o altos abiertos (RT-11.20) (BTT, p. 24) \\
Perfil operacional & Dispositivos de producción, uno de cada tipo; personas del CLIENTE convocadas & Ver sección 7.3 \\
Migración & Perfilado y saneamiento cerrados (3.7.1); mapeo del SD5, Anexo 5-K & Conciliación sin diferencias no explicadas en cada ensayo; tiempos medidos compatibles con la ventana y con la reversión (SD5, 5.3.5) \\
Despliegue sin interrupción & Estrategia canario configurada & Despliegue y reversión sin interrupción del servicio y sin intervención manual (RT-04.06, RT-04.07) (BTT, p. 11) \\
\end{tablasimple}
\fuentesd{Fuente: elaboración a partir de Distribuidora Puelche S.A. (2026a, arts. 17.3 y 24, pp. 13, 17); Distribuidora Puelche S.A. (2026b, sección 20.1, pp. 34--35); Distribuidora Puelche S.A. (2026c, cap. 15); LafroX (2026q, Anexos 9, anexos 9.A y 9.B); LafroX (2026f, T-17, sección 4).}

\subsection*{6.2 Suspensión y reanudación}
\addcontentsline{toc}{subsection}{6.2 Suspensión y reanudación}

Una campaña se suspende cuando aparece un defecto crítico que impide seguir ejecutando casos, cuando el ambiente pierde su equivalencia con Producción o cuando los datos dejan de estar en el estado declarado. Se reanuda con el defecto corregido y verificado en QA, con el ambiente reconstruido y con los datos restituidos. Los casos ya aprobados que dependen del componente corregido se ejecutan de nuevo.

\section*{7. Especificación de las pruebas que habilitan cada entrada a Producción}
\addcontentsline{toc}{section}{7. Especificación de las pruebas que habilitan cada entrada a Producción}

Las pruebas siguientes comprueban carga, fallas, operación de terreno y continuidad híbrida antes de habilitar la ola aplicable.

\subsection*{7.1 Carga y estrés}
\addcontentsline{toc}{subsection}{7.1 Carga y estrés}

La carga se deriva del dimensionamiento del SD4, sección 4.2.6, y no se elige:

\begin{tablasimple}{>{\hsize=1.899\hsize\linewidth=\hsize}Y >{\hsize=0.7\hsize\linewidth=\hsize}Y >{\hsize=0.7\hsize\linewidth=\hsize}Y >{\hsize=0.701\hsize\linewidth=\hsize}Y}
  {\cab{Magnitud} & \cab{Peak declarado (SD4)} & \cab{Factor} & \cab{Carga de prueba}}
TPS máximo de septiembre (12:00) & 14,66 TPS & 1,5 (BA, Art. 24, p. 17; RT-09.06) & 21,99 TPS \\
TPS máximo de la ventana de 05:30 a 07:00 & 6,94 TPS & 1,5 & 10,41 TPS, sostenidos 90 minutos \\
Personas usuarias concurrentes (cota extrema) & 775,33 & 1,5 & 1.163 \\
Dispositivos de terreno simultáneos & 158 & 1,5 & 237 \\
Escalón de crecimiento (RT-09.03) (BTT, p. 21) & 14,66 TPS & 3 & 43,98 TPS \\
Escalón de quiebre & 14,66 TPS & 4 & 58,64 TPS \\
\end{tablasimple}
\fuentesd{Fuente: elaboración a partir de Distribuidora Puelche S.A. (2026a, art. 24, p. 17); Distribuidora Puelche S.A. (2026b, RT-09.03 y RT-09.06, p. 21); LafroX (2026m, SD4, secciones 4.2.6.3, 4.2.6.4 y 4.2.6.7); LafroX (2026l, Anexos 4, secciones 4-W.2, 4-W.3 y 4-W.6).}

\textbf{Umbrales de salida (p95, SD4, Tabla 21):} confirmación de línea de preparación ≤ 1 s; registro de entrega ≤ 2 s; línea de preventa ≤ 1,5 s; consulta de stock y crédito ≤ 2 s; transacción crítica de terreno ≤ 3 s; carga inicial del portal ≤ 2 s; navegación ≤ 1 s; API de consulta ≤ 500 ms; API de escritura ≤ 800 ms; búsqueda compuesta ≤ 3 s y búsqueda global ≤ 2 s; informe estándar ≤ 30 s.

\textbf{Ventana de despacho.} La prueba de 10,41 TPS concentra la emisión de guías en los 90 minutos. Su criterio de salida incluye 0 minutos de indisponibilidad (Caso, RT-10.05).

\textbf{Informe (RT-09.07):} curva de tiempo de respuesta frente a carga, punto de saturación, consumo de recursos por componente, y comportamiento durante y después del peak. El informe identifica el primer cuello de botella declarado en el SD4 (RT-09.05) y compara su comportamiento observado con el previsto (BTT, p. 21).

\subsection*{7.2 Resiliencia}
\addcontentsline{toc}{subsection}{7.2 Resiliencia}

La inyección reproduce fallas de nube, enlace y equipamiento local. El resultado se compara con la continuidad y los tiempos comprometidos, sin sustituir una falla por una excepción.

\begin{tablasimple}{>{\hsize=0.808\hsize\linewidth=\hsize}Y >{\hsize=0.771\hsize\linewidth=\hsize}Y >{\hsize=1.421\hsize\linewidth=\hsize}Y}
  {\cab{Escenario (RT-10.07)} (BTT, p. 22) & \cab{Forma de inyección} & \cab{Resultado esperado}}
Caída de instancia & AWS FIS detiene tareas ECS Fargate & Reemplazo automático sin pérdida \\
Caída de zona de disponibilidad & AWS FIS sobre una zona & Continuidad en las zonas restantes \\
Dependencia externa sin respuesta, con error o lenta & Simulador de la integración afectada & Comportamiento declarado por integración (RT-10.08) (BTT, p. 22) \\
Latencia elevada & Inyección de latencia de red & Umbrales de la sección 7.1 o degradación controlada \\
Saturación de disco & Llenado controlado del volumen & Alerta anticipada y continuidad \\
Corte de enlace del centro de distribución por 24 h & Desconexión del sitio emulado & Operación local continua; sincronización ≤ 2 h tras reconectar, con conflictos de stock resueltos de forma determinista (RT-03.10, RT-03.13) \\
Pérdida de fibra y LTE de un centro de distribución & Corte de ambos enlaces en el sitio emulado & Conmutación a Starlink en espera caliente; la emisión de guías por el ERP continúa (SD4, sección 4.1.3.2; AL-DTE-01) \\
Caída de un gateway de frío de Talca & Detención de uno de los dos gateways con Greengrass & El otro gateway lee ambas cámaras; el bloqueo térmico local actúa en ≤ 5 s y la alerta llega a M9 (SD4, ADR-22) \\
\end{tablasimple}
\fuentesd{Fuente: elaboración a partir de Distribuidora Puelche S.A. (2026b, RT-10.07, p. 22); LafroX (2026m, SD4, sección 4.2.4.3); LafroX (2026l, Anexos 4, anexo 4-V).}

Los dos últimos escenarios prueban la parte on-premise de la arquitectura híbrida, que el RT-10.07 no enumera pero que el SD4 declara crítica para el despacho y para la cadena de frío (BTT, p. 22).

\subsection*{7.3 Perfil operacional de terreno}
\addcontentsline{toc}{subsection}{7.3 Perfil operacional de terreno}

Las pruebas usan dispositivos reales y personas del perfil correspondiente. Una prueba de laboratorio no reemplaza las condiciones físicas del turno y la descarga.

\begin{tablasimple}{>{\hsize=0.637\hsize\linewidth=\hsize}Y >{\hsize=1.071\hsize\linewidth=\hsize}Y >{\hsize=1.292\hsize\linewidth=\hsize}Y}
  {\cab{Condición (Caso, cap. 15)} & \cab{Prueba} & \cab{Criterio de salida}}
Turno completo de 14 h sin cobertura & Jornada completa de un repartidor rural, incluida la ruta de 34 clientes (SD4, sección 4.2.6.7) & 0 registros perdidos o duplicados; sincronización ≤ 10 min tras reconectar \\
\ensuremath{-}22 °C con guantes térmicos & Preparación y recepción en cámara de congelado & Transacción completada con guantes; indicadores RT-13.04 del perfil dentro de umbral (BTT, p. 26) \\
Lluvia & Entrega a la intemperie bajo lluvia & Transacción completada; dispositivo operativo \\
Sobre 34 °C & Entrega a la intemperie con calor & Transacción completada; dispositivo operativo \\
Sol directo & Lectura de la pantalla en la puerta del local & Lectura sin error de la persona usuaria \\
Una sola mano durante la descarga & Registro de entrega & Transacción completada con una mano \\
\end{tablasimple}
\fuentesd{Fuente: elaboración a partir de Distribuidora Puelche S.A. (2026c, cap. 15); Distribuidora Puelche S.A. (2026b, RT-13.04, p. 26); LafroX (2026m, SD4, sección 4.2.6.7); LafroX (2026l, Anexos 4, sección 4-W.6).}

Cada perfil mide los cuatro indicadores del RT-13.04 con los valores del SD3, Tabla 3.A.4a (métrica 9A-18): tasa de error ≤ 5 \% en todos los perfiles y aprendizaje ≤ 2 h en bodega, ≤ 4 h para el preventista y ≤ 30 min para el conductor (BTT, p. 26).

\subsection*{7.4 Pruebas de los resultados del caso}
\addcontentsline{toc}{subsection}{7.4 Pruebas de los resultados del caso}

Las pruebas técnicas no aceptan por sí solas los resultados del Caso, cap. 18. Algunos se ensayan antes de la marcha blanca para llegar a ella sin sorpresas:

\begin{tablasimple}{>{\hsize=1.055\hsize\linewidth=\hsize}Y >{\hsize=1.378\hsize\linewidth=\hsize}Y >{\hsize=0.567\hsize\linewidth=\hsize}Y}
  {\cab{Resultado} & \cab{Ensayo previo} & \cab{Cuándo}}
R18-01: clientes afectados por un lote en menos de 2 h & Simulacro de retiro sanitario sobre datos sintéticos con lotes trazados & Aceptación de la Etapa 1 (3.8.2) \\
R18-03: registro continuo de temperatura & Series con excursiones controladas; verificación del bloqueo del despacho ante una excursión fuera de rango & Integración (3.8.1) y aceptación (3.8.2) \\
R18-06: ningún pedido perdido ni duplicado por falta de señal & Jornada de 14 h sin señal (sección 7.3) & Perfil operacional (3.8.3) \\
R18-07: ruta del día siguiente en menos de 20 minutos & Generación de la ruta con el volumen de un día de peak & Carga (3.8.4) \\
\end{tablasimple}
\fuentesd{Fuente: elaboración a partir de Distribuidora Puelche S.A. (2026c, cap. 18); LafroX (2026h, SD13, sección 13.2); LafroX (2026f, T-17, sección 5).}

Los dieciséis resultados se aceptan con el protocolo del Formulario T-17, en el momento que fija el SD7, Tabla 7.11.

\subsection*{7.5 Protocolos de aceptación de la arquitectura híbrida}
\addcontentsline{toc}{subsection}{7.5 Protocolos de aceptación de la arquitectura híbrida}

El SD4, Anexo 4-V, define nueve protocolos que verifican la continuidad de la arquitectura en la nube y en los cinco sitios con servicio local. Se ejecutan en Preproducción antes de la ola aplicable, dentro de las campañas de este plan:

\begin{tablasimple}{Y Y Y}
  {\cab{Protocolo} & \cab{Qué verifica (SD4, Anexo 4-V)} & \cab{Campaña de este plan}}
AL-DTE-01 & 96 guías por las rutas de Talca y de los demás sitios, con fallas y cambio de carga; emisión por fibra, LTE o Starlink & Resiliencia (3.8.4; 3.9.3) \\
AL-POD-01 & Entrega total, parcial y con receptor sustituto sin señal & Perfil operacional (3.8.3) \\
AL-SLA-01 & Ventanas y cuotas del ERP, del SII y de las contrapartes & Integración (3.8.1; 3.9.1) \\
AL-OFF-01 & Aislamiento de cada sitio por 24 horas, con relevos y despacho con guía preemitida & Resiliencia (3.8.4; 3.9.3) \\
AL-CLI-01 & Pedidos sin señal en el Portal de Clientes instalado y reconexión & Integración (3.8.1; 3.9.1) \\
AL-DR-01 & Corte de los caminos y recuperación del WMS de Talca en la nube & Recuperación ante desastres (3.8.5; 3.9.4) \\
AL-PERF-01 & Carga de 3× con mamparos por perfil & Carga y estrés (3.8.4; 3.9.3) \\
AL-STOCK-01 & Última unidad pedida desde preventa, portal y EDI, con cortes y acuses perdidos & Integración (3.8.1; 3.9.1) \\
AL-ACT-01 & Cada actor con una operación permitida y una denegada & Integración (3.8.1; 3.9.1) \\
\end{tablasimple}
\fuentesd{Fuente: elaboración a partir de LafroX (2026l, Anexos 4, anexo 4-V, protocolos AL-DTE-01 a AL-ACT-01); LafroX (2026d, T-14, paquetes indicados en cada fila).}

Cada protocolo deja su acta con la carga, la falla inyectada, el resultado observado y los responsables, y esa acta forma parte del expediente del H5 o del H10 (T-17, sección 3.1).

\section*{8. Automatización}
\addcontentsline{toc}{section}{8. Automatización}

Se automatizan los flujos repetibles y se conserva la intervención humana donde el juicio, la condición física o la independencia de la prueba lo requieren.

\subsection*{8.1 Criterio}
\addcontentsline{toc}{subsection}{8.1 Criterio}

La automatización es la regla. Una prueba es manual solo si exige el juicio de una persona usuaria (aceptación, usabilidad, accesibilidad manual), una condición física (perfil operacional) o un tercero independiente (seguridad ofensiva).

\subsection*{8.2 Tamaño estimado del catálogo}
\addcontentsline{toc}{subsection}{8.2 Tamaño estimado del catálogo}

El catálogo se dimensiona por riesgo, con la prioridad de cada requisito en el SD3, Anexos. Los casos por requisito (cinco, tres y dos según la prioridad; uno por RT; siete por integración) forman la regla de dimensionamiento por riesgo de LafroX, que el plan del H1 contrasta con el diseño real de los casos:

\begin{tablasimple}{>{\hsize=1.719\hsize\linewidth=\hsize}Y >{\hsize=0.76\hsize\linewidth=\hsize}Y >{\hsize=0.76\hsize\linewidth=\hsize}Y >{\hsize=0.761\hsize\linewidth=\hsize}Y}
  {\cab{Grupo} & \cab{Cantidad} & \cab{Casos por requisito} & \cab{Casos}}
RF y RNF de prioridad crítica & 56 & 5 & 280 \\
RF y RNF de prioridad alta & 163 & 3 & 489 \\
RF y RNF de prioridad media & 43 & 2 & 86 \\
RT en estado «Cumple» o «Cumple parcialmente» en el T-12 & 356 & 1 & 356 \\
Contratos de integración & 15 & 7 & 105 \\
\textbf{Total} &  &  & \textbf{1.316} \\
\end{tablasimple}
\fuentesd{Fuente: elaboración a partir de LafroX (2026k, SD3, catálogo de RF y RNF); LafroX (2026b, T-12, catálogo de RT ofertados); LafroX (2026m, SD4, sección 4.1); las reglas de dimensionamiento 5/3/2, un caso por RT y siete por integración de esta sección 8.2.}

Los 262 RF y RNF con prioridad propia son los del catálogo del SD3. Las nueve filas RF y RNF restantes del T-12 (271 en total) se verifican con la fila que las contiene y no suman casos. Los 18 RT en estado «No cumple» no tienen casos; quedan en la matriz de trazabilidad como brecha visible.

\subsection*{8.3 Duración de la regresión}
\addcontentsline{toc}{subsection}{8.3 Duración de la regresión}

El cálculo usa tres parámetros de diseño del plan: 80 \% de los casos automatizados, 30 segundos por caso automático y 15 minutos por caso manual. Con ellos: 0,8 × 1.316 \ensuremath{\approx} 1.053 casos automáticos. A 30 segundos por caso, la regresión completa toma 1.053 × 0,5 / 60 \ensuremath{\approx} 8,8 horas en un ejecutor y \ensuremath{\approx} 4,4 horas en dos ejecutores paralelos. Cabe en una ejecución nocturna antes de cada promoción a Preproducción.

Los 263 casos manuales, a 15 minutos cada uno, requieren 263 × 15 / 60 \ensuremath{\approx} 65,8 HH por ciclo. Se planifican con los evaluadores en los meses de certificación.

\section*{9. Calendario}
\addcontentsline{toc}{section}{9. Calendario}

El calendario vincula las campañas con los hitos y distingue implementación, marchas blancas y Operación. Las subventanas de ejecución se controlan sin cambiar los meses contractuales.

\subsection*{9.1 Implementación}
\addcontentsline{toc}{subsection}{9.1 Implementación}

Las fechas delimitan el trabajo de cada paquete, incluida su preparación e informe; no significan que todas las intervenciones sean simultáneas.

\begin{tablasimple}{>{\hsize=2.052\hsize\linewidth=\hsize}Y >{\hsize=0.79\hsize\linewidth=\hsize}Y >{\hsize=0.79\hsize\linewidth=\hsize}Y >{\hsize=0.79\hsize\linewidth=\hsize}Y >{\hsize=0.79\hsize\linewidth=\hsize}Y >{\hsize=0.788\hsize\linewidth=\hsize}Y}
  {\cab{Campaña} & \cab{Paquete} & \cab{Instancia o ambiente} & \cab{Inicio} & \cab{Término} & \cab{HH (T-15)}}
Plan de calidad y de pruebas; revisión preliminar el 05-03-2027 & 1.5.1 & — & 01-03-2027 & 09-03-2027 & 80 \\
Investigación con personas usuarias & 2.6.1 & Puestos de trabajo & 01-03-2027 & 29-03-2027 & 240 \\
Matriz de trazabilidad (H1) & 1.2.4 & — & 01-03-2027 & 09-03-2027 & 80 \\
Prototipos y pruebas de usabilidad & 2.6.2 & Prototipos & 01-04-2027 & 29-04-2027 & 240 \\
Indicadores de usabilidad & 2.6.3 & — & 03-05-2027 & 31-05-2027 & 240 \\
Puertas automáticas activas & 1.5.2 & Cadena de integración & Mes 6 & Mes 20 & 80 \\
Ensayo de migración 1 & 3.7.2 & PRE-M & 04-10-2027 & 03-11-2027 & (en 480 HH del paquete) \\
Integración, regresión e idempotencia E1 (H4) & 3.8.1 & QA & 01-10-2027 & 20-10-2027 & 160 \\
Aceptación y accesibilidad E1 & 3.8.2 & PRE-A & 21-10-2027 & 17-11-2027 & 320 \\
Perfil operacional E1 & 3.8.3 & Terreno y cámara & 21-10-2027 & 17-11-2027 & 320 \\
Carga, estrés y resiliencia E1 & 3.8.4 & PRE-C & 21-10-2027 & 09-11-2027 & 320 \\
Recuperación ante desastres E1 & 3.8.5 & Recuperación & 21-10-2027 & 09-11-2027 & 320 \\
Seguridad ofensiva E1 & 3.8.6 & PRE-S & 21-10-2027 & 09-11-2027 & 320 \\
Despliegue sin interrupción E1 & 3.8.8 & PRE-C & 21-10-2027 & 09-11-2027 & 320 \\
Acta de certificación E1 (H5) & 3.8.7 & — & 18-11-2027 & 29-11-2027 & 160 \\
Ensayo de migración 2 & 3.7.3; 3.7.4 & PRE-M & 01-12-2027 & 31-12-2027 & (en las HH de los paquetes) \\
Acta de conciliación y corte & 3.7.5 & — & 03-01-2028 & 27-01-2028 & 480 \\
Integración y regresión E2 sin afectar E1 (H9) & 3.9.1 & QA & 01-05-2028 & 18-05-2028 & 160 \\
Carga con ambas etapas y resiliencia & 3.9.3 & PRE-C & 19-05-2028 & 26-05-2028 & 320 \\
Recuperación ante desastres con ambas etapas & 3.9.4 & Recuperación & 19-05-2028 & 26-05-2028 & 320 \\
Seguridad ofensiva E2 & 3.9.5 & PRE-S & 19-05-2028 & 26-05-2028 & 160 \\
Despliegue sin interrupción E2 & 3.9.7 & PRE-C & 19-05-2028 & 26-05-2028 & 160 \\
Aceptación y accesibilidad E2 & 3.9.2 & PRE-A & 27-05-2028 & 31-05-2028 & 320 \\
Acta de certificación E2 (H10) & 3.9.6 & — & 01-06-2028 & 12-06-2028 & 160 \\
\end{tablasimple}
\fuentesd{Fuente: elaboración a partir de LafroX (2026d, T-14, paquetes indicados en cada fila); LafroX (2026e, T-15, sección 6.1); LafroX (2026f, T-17, sección 4.1).}

En E2, 3.9.1/3.9.2 incluyen regresión de las cinco condiciones físicas y de desconexión sobre funciones/dispositivos modificados, con evidencia de cobertura para condiciones no modificadas; no se omite el perfil por ausencia de un paquete nuevo.

La campaña de aceptación de la Etapa 2 es breve porque cubre solo el alcance incremental. Los flujos de la Etapa 1 se verifican con la regresión de 3.9.1, que corre antes.

La ventana 21-10 a 09-11-2027 de 3.8.5 comprende preparación no invasiva del 21 al 25 de octubre, ejecución/conmutación y reversión del 26 al 29 de octubre, y análisis/informe hasta el 9 de noviembre. No se interviene el sitio el 1 de noviembre (feriado) ni los días 2, 3 y 4 (primeros tres hábiles efectivos); el análisis no modifica el sitio. La autorización precisa además horario fuera de despacho y reversión segura.

Ninguna ejecución invasiva de campaña que intervenga Producción o el sitio de recuperación cae en un congelamiento (septiembre y diciembre) ni en los tres primeros días hábiles del mes (Caso, RT-10.05). Las campañas en Preproducción pueden hacerlo, porque no tocan Producción.

\subsection*{9.2 Marchas blancas y soporte puente}
\addcontentsline{toc}{subsection}{9.2 Marchas blancas y soporte puente}

Durante la operación supervisada se miden indicadores y se prueban respaldos. Cada corrección repite la verificación del alcance afectado antes de promoverse.

\begin{tablasimple}{>{\hsize=0.56\hsize\linewidth=\hsize}Y >{\hsize=2.062\hsize\linewidth=\hsize}Y >{\hsize=0.378\hsize\linewidth=\hsize}Y}
  {\cab{Período} & \cab{Prueba} & \cab{Frecuencia}}
Meses 13 a 15 (E1) y 19 a 20 (E2) & Indicadores diarios del Art. 17.3 (SD7, Tabla 7.9) y conciliación con el registro vigente (BA, p. 13) & Diaria \\
Meses 13 a 20 & Restauración de respaldos sobre una muestra representativa, con tiempo medido (RT-07.12), por SRE: 4.2.1 (13–15), 4.2.2 (16–18) y 4.3.1 (19–20); ocho informes distintos de los 36 de 8.1.6 (BTT, p. 18) & Mensual \\
Cada corrección & Puertas G0 a G4 y G6, regresión completa, DAST y carga de regresión; repetición de la prueba de certificación afectada & Por corrección \\
\end{tablasimple}
\fuentesd{Fuente: elaboración a partir de Distribuidora Puelche S.A. (2026a, art. 17.3, p. 13); LafroX (2026p, SD7, Tabla 7.9); LafroX (2026d, T-14, cuenta 4).}

Cada restauración previa tiene 16 HH SRE (4 preparación, 4 ejecución, 4 validación y 4 informe); se agregan 48 HH a 4.2.1, 48 a 4.2.2 y 32 a 4.3.1 como componentes separados de la cobertura NOC/mesa. Las siete conmutaciones operacionales requieren 560 HH y la recuperación inicial 320 HH según el desglose por actividades/roles del T-15; el mismo formulario justifica 320/160 HH de intrusión E1/E2.

\subsection*{9.3 Operación (meses 21 a 56)}
\addcontentsline{toc}{subsection}{9.3 Operación (meses 21 a 56)}

Las pruebas periódicas verifican la continuidad del servicio durante los 36 meses. Su frecuencia y sus informes son distintos de las certificaciones iniciales.

\begin{tablasimple}{>{\hsize=1.447\hsize\linewidth=\hsize}Y >{\hsize=1.41\hsize\linewidth=\hsize}Y >{\hsize=0.655\hsize\linewidth=\hsize}Y >{\hsize=1.334\hsize\linewidth=\hsize}Y >{\hsize=0.577\hsize\linewidth=\hsize}Y >{\hsize=0.577\hsize\linewidth=\hsize}Y}
  {\cab{Prueba} & \cab{Exigencia} & \cab{Meses} & \cab{Fechas} & \cab{Cantidad} & \cab{Paquete}}
Recuperación ante desastres con conmutación real & Al menos dos veces al año (RT-07.07) (BTT, p. 17) & 21, 27, 33, 39, 45, 51, 54 & oct-28, abr-29, oct-29, abr-30, oct-30, abr-31, jul-31 & 7 & 8.1.3 \\
Resiliencia por inyección de fallas & Al menos una vez por semestre (RT-10.07) (BTT, p. 22) & 26, 32, 38, 44, 50, 55 & mar-29, sep-29, mar-30, sep-30, mar-31, ago-31 & 6 & 8.2.2 \\
Intrusión por tercero independiente & Anual (RT-11.20) (BTT, p. 24) & 28, 40, 52 & may-29, may-30, may-31 & 3 & 8.2.2 \\
Restauración de respaldos & Al menos mensual (RT-07.12) (BTT, p. 18) & 21 a 56 & Mensual & 36 & 8.1.6 \\
Regresión de INN-02 & Escenarios reproducidos de incidentes & 21 a 56 & Trimestral (informe) & 12 & 8.3.1 \\
\end{tablasimple}
\fuentesd{Fuente: elaboración a partir de Distribuidora Puelche S.A. (2026b, RT-07.07, RT-07.12, RT-10.07 y RT-11.20, pp. 17--18, 22, 24); Distribuidora Puelche S.A. (2026c, RT-10.05); LafroX (2026d, T-14, cuenta 8); LafroX (2026e, T-15, sección 6.1).}

Las resiliencias de septiembre (meses 32 y 44) se ejecutan en Preproducción, porque septiembre está congelado. Cada prueba periódica produce un informe con el resultado medido, la comparación con su umbral y el plan de corrección de las brechas (RT-07.07) (BTT, p. 17).

\section*{10. Defectos}
\addcontentsline{toc}{section}{10. Defectos}

La severidad se asigna según el servicio afectado y determina el bloqueo de la promoción. El registro conserva requisito, versión, prueba y evidencia de cierre.

\begin{tablasimple}{>{\hsize=0.338\hsize\linewidth=\hsize}Y >{\hsize=1.846\hsize\linewidth=\hsize}Y >{\hsize=0.816\hsize\linewidth=\hsize}Y}
  {\cab{Severidad} & \cab{Definición (alineada con el SD4, sección 4.2.4.3)} & \cab{Efecto en las pruebas}}
Crítica & Afecta un servicio de nivel crítico: preparación y despacho con emisión de guía en la ventana de 05:30 a 07:00, bloqueo por excursión térmica o identidad de bodega & Detiene la promoción; puede suspender la campaña \\
Alta & Afecta un servicio de nivel alto: preventa y consulta de stock y crédito, entrega y su evidencia, rutas, documentos tributarios, ERP, cobranza, registro de temperatura & Detiene la promoción \\
Media & Afecta un servicio de nivel medio con alternativa operativa & Se registra y se planifica \\
Baja & Afecta la reportería sin impedir operar & Se registra y se planifica \\
\end{tablasimple}
\fuentesd{Fuente: elaboración a partir de Distribuidora Puelche S.A. (2026a, art. 78.2, p. 40); LafroX (2026m, SD4, niveles de servicio); los efectos sobre las campañas definidos en esta sección 10.}

Cada defecto se registra en GitLab con el caso de prueba, el requisito, el ambiente, la versión y la evidencia. Se cierra solo cuando la prueba que lo detectó pasa y queda incorporada a la regresión. Los hallazgos de seguridad conservan su severidad propia y sus plazos de corrección de 7, 15 y 30 días corridos para severidad crítica, alta y media (BA, Art. 21.1, p. 15; SD4, sección 4.2.4.1.4).

\section*{11. Roles}
\addcontentsline{toc}{section}{11. Roles}

Las responsabilidades siguientes separan el diseño y ejecución de las pruebas de la aceptación contractual que corresponde al CLIENTE.

\begin{tablasimple}{>{\hsize=0.528\hsize\linewidth=\hsize}Y >{\hsize=1.156\hsize\linewidth=\hsize}Y >{\hsize=1.316\hsize\linewidth=\hsize}Y}
  {\cab{Rol} & \cab{Persona o equipo} & \cab{Responsabilidad en este plan}}
Líder de Calidad y Pruebas & Maximiliano Miño & Dueño del plan; aprueba el inicio y el cierre de cada campaña; veto sobre la promoción \\
Equipo de pruebas (CAL) & Analistas de pruebas de LafroX, con refuerzo de evaluadores subcontratados en los meses de certificación (SD6, sección 6.1.3) & Diseño, automatización y ejecución \\
Seguridad (SEG) & Encargado de Seguridad, Álvaro Catalán, y tercero independiente & Análisis de seguridad y pruebas de intrusión \\
Operación y confiabilidad (SRE) & Líder de Operación, Guillermo Castillo & Resiliencia, recuperación, restauración, despliegue sin interrupción \\
Datos (DAT) & Líder de Datos, Leandro Chamorro & Ensayos de migración y conciliación \\
Líder Funcional & Patricio Henríquez & Coordina la aceptación con las personas del CLIENTE \\
Contraparte Técnica del CLIENTE & Designada por Distribuidora Puelche & Firma los casos de aceptación y las actas (T-17) \\
\end{tablasimple}
\fuentesd{Fuente: elaboración a partir de LafroX (2026i, SD1, sección 1.5); LafroX (2026o, SD6, organización de especialistas); LafroX (2026e, T-15, asignación de recursos).}

\section*{12. Informes de prueba}
\addcontentsline{toc}{section}{12. Informes de prueba}

Los informes conservan la comparación entre resultado observado y criterio de salida, y se entregan al destinatario indicado con la frecuencia del plan.

\begin{tablasimple}{>{\hsize=0.857\hsize\linewidth=\hsize}Y >{\hsize=1.96\hsize\linewidth=\hsize}Y >{\hsize=0.741\hsize\linewidth=\hsize}Y >{\hsize=0.442\hsize\linewidth=\hsize}Y}
  {\cab{Informe} & \cab{Contenido mínimo} & \cab{Destinatario} & \cab{Frecuencia}}
Informe de ejecución de la cadena de integración & Resultados por requisito, cobertura, métricas de código & Equipo y Líder de Calidad & Cada ejecución \\
Informe de campaña & Casos ejecutados, aprobados y fallidos; defectos por severidad; desviaciones del plan; conclusión de salida & Líder de Calidad y Contraparte Técnica & Al cierre de cada campaña \\
Informe de carga & Contenido del RT-09.07 (sección 7.1) (BTT, p. 21) & Contraparte Técnica & Por campaña de carga \\
Informe de recuperación ante desastres & RTO y RPO medidos; brechas y su plan de corrección (RT-07.07) (BTT, p. 17) & Contraparte Técnica & Por prueba \\
Informe íntegro de intrusión & Hallazgos, severidad, plan de remediación con plazos (RT-11.20) (BTT, p. 24) & CLIENTE & Por prueba \\
Informe mensual de calidad & Métricas de proceso, defectos escapados, estado de la matriz de trazabilidad & Comité de Calidad y Contraparte Técnica & Mensual \\
\end{tablasimple}
\fuentesd{Fuente: elaboración a partir de Distribuidora Puelche S.A. (2026b, RT-09.07, RT-07.07 y RT-11.20, pp. 17, 21, 24); ISO, IEC \& IEEE (2021c, documentación de pruebas); LafroX (2026f, T-17, sección 3).}

\section*{Referencias}
\addcontentsline{toc}{section}{Referencias}

Las fuentes siguientes sustentan los criterios citados en este documento; sus referencias se leen junto con el artículo, requisito o sección indicado en el lugar de uso.

\begin{itemize}[before=\raggedright,itemsep=0pt,parsep=0pt,topsep=2pt]
  \item Distribuidora Puelche S.A. (2026a). \emph{Bases Administrativas, Licitación TFEP-01/2026}.
  \item Distribuidora Puelche S.A. (2026b). \emph{Bases Técnicas Transversales, Licitación TFEP-01/2026}.
  \item Distribuidora Puelche S.A. (2026c). \emph{Bases Técnicas del Caso Logística}.
  \item International Organization for Standardization, International Electrotechnical Commission, \& Institute of Electrical and Electronics Engineers. (2021a). \emph{ISO/IEC/IEEE 29119-2:2021. Software and systems engineering — Software testing — Part 2: Test processes}.
  \item International Organization for Standardization, International Electrotechnical Commission, \& Institute of Electrical and Electronics Engineers. (2021c). \emph{ISO/IEC/IEEE 29119-3:2021. Software and systems engineering — Software testing — Part 3: Test documentation}.
  \item International Organization for Standardization, International Electrotechnical Commission, \& Institute of Electrical and Electronics Engineers. (2021b). \emph{ISO/IEC/IEEE 29119-4:2021. Software and systems engineering — Software testing — Part 4: Test techniques}.
  \item International Organization for Standardization, International Electrotechnical Commission, \& Institute of Electrical and Electronics Engineers. (2022). \emph{ISO/IEC/IEEE 29119-1:2022. Software and systems engineering — Software testing — Part 1: General concepts}.
  \item Ley N.º 21.719. (2024). \emph{Regula la protección y el tratamiento de los datos personales y crea la Agencia de Protección de Datos Personales}. Diario Oficial de la República de Chile.
  \item LafroX. (2026a). \emph{Formulario T-10: Metodología de desarrollo de software}. Propuesta técnica, Licitación TFEP-01/2026.
  \item LafroX. (2026b). \emph{Formulario T-12: Matriz de trazabilidad}. Propuesta técnica, Licitación TFEP-01/2026.
  \item LafroX. (2026c). \emph{Formulario T-13: Plan de pruebas y validación}. Propuesta técnica, Licitación TFEP-01/2026.
  \item LafroX. (2026d). \emph{Formulario T-14: Estructura de desglose del trabajo}. Propuesta técnica, Licitación TFEP-01/2026.
  \item LafroX. (2026e). \emph{Formulario T-15: Cronograma y asignación de recursos}. Propuesta técnica, Licitación TFEP-01/2026.
  \item LafroX. (2026f). \emph{Formulario T-17: Protocolo de aceptación}. Propuesta técnica, Licitación TFEP-01/2026.
  \item LafroX. (2026h). \emph{Subdocumento 13: Innovaciones}. Propuesta técnica, Licitación TFEP-01/2026.
  \item LafroX. (2026i). \emph{Subdocumento 1: Presentación de la empresa}. Propuesta técnica, Licitación TFEP-01/2026.
  \item LafroX. (2026k). \emph{Subdocumento 3: Especificación de requerimientos}. Propuesta técnica, Licitación TFEP-01/2026.
  \item LafroX. (2026l). \emph{Subdocumento 4: Anexos}. Propuesta técnica, Licitación TFEP-01/2026.
  \item LafroX. (2026m). \emph{Subdocumento 4: Arquitectura de la solución}. Propuesta técnica, Licitación TFEP-01/2026.
  \item LafroX. (2026n). \emph{Subdocumento 5: Modelo y gestión de datos}. Propuesta técnica, Licitación TFEP-01/2026.
  \item LafroX. (2026o). \emph{Subdocumento 6: Metodología de desarrollo}. Propuesta técnica, Licitación TFEP-01/2026.
  \item LafroX. (2026p). \emph{Subdocumento 7: Plan de trabajo}. Propuesta técnica, Licitación TFEP-01/2026.
  \item LafroX. (2026q). \emph{Subdocumento 9: Anexos 9.A a 9.D}. Propuesta técnica, Licitación TFEP-01/2026.
  \item LafroX. (2026r). \emph{Subdocumento 9: Plan de calidad}. Propuesta técnica, Licitación TFEP-01/2026.
\end{itemize}
